\subsection{Methodik}
\label{method_oluanw}

\subsubsection{Grundlegendes Vorgehen}
\label{vorgehen_oluanw}
Off-Label-Anwendungen sind definiert als Kombination einer Off-Label-Art und
eines Wirkstoffs. Wir wollen die Abhängigkeiten von häufig auftretenden
Off-Label-Anwendungen zwsischen den verschiedenen Diagnosen und Einrichtungen
deskriptiv untersuchen. Hierbei berechnen wir den Anteil jeder
Off-Label-Anwendung an allen Off-Label-Verordnungen, in der jeweiligen
Subgruppe, die durch die entsprechende Variable definiert ist. Als
Off-Label-Verordnungen bezeichnen wir jede Verordnungen die ganz oder teilweise
Off-Label waren. Um die Off-Label-Anwendungen besser vergleichbar zu machen,
ordnen wir sie entsprechend ihres $\chi^2$-Koeffizienten \citep{Fahrmeir2016}.
Aufgrund der begrenzten Datenmenge und der Vielzahl von Off-Label-Anwendungen,
die nur selten auftreten, haben wir uns entschlossen, nur jene
Off-Label-Anwendungen zu berücksichtigen, die mindestens 20 Mal vorkamen. 

\subsubsection{\texorpdfstring{$\chi^2$}{X²}-Koeffizient}
\label{chi_quadrat_oluanw}
Der $\chi^2$-Koeffizient ist ein Maß für die Stärke der Abhängigkeit zwischen
zwei Variablen. Für jede Kombination aus einer Off-Label-Anwendung $Anwendung^a,
a \in \{1, \dots, A\}$ und einer Ausprägung $v_k, k = 1,\dots,K$ der Merkmale $V
\in \{Diagnose, Einrichtung2\}$ wird eine Häufigkeitstabelle
(\ref{tab:frequenz_oluanw}) erstellt, die die beobachtete relative Anzahl der
Fälle für jede Kombination von $Anwendung^a$ und $v_k$ enthält. Die Werte
$f_{k1}, k = 1,\dots, K$ geben die relative Häufigkeit der beobachteten
Off-Label-Verordnungen mit Off-Label-Anwendung $Anwendung^a$ in Einrichtung bzw.
Diagnose $v_k$, relativ zu der Gesamtanzahl der Off-Label-Verordnungen, an.
$f_{k\cdot}$ ist somit die Anteil aller beobachteten Off-Label-Verordnungen mit
Off-Label-Anwendung $Anwendung^a$. Die Werte $f_{k2}$ sind analog definiert für
beobachteten Off-Label-Verordnungen mit einer anderen Off-Label-Anwendung als
$Anwendung^a$. Die Gesamtanzahl an Off-Label-Verordnungen bezeichnen wir mit
$n$.

\begin{equation}
\begin{tabular}{c|c|c|c}
    \multicolumn{1}{c}{} & \multicolumn{1}{c}{$Anwendung^a$} & \multicolumn{1}{c}{$\overline{Anwendung^a}$} & \\
    \cline{2-3}
    $v_1$ & $f_{11}$ & $f_{12}$ & $f_{1\cdot}$ \\
    \vdots & \vdots & \vdots & \vdots \\
    $v_k$ & $f_{k1}$ & $f_{k2}$ & $f_{k\cdot}$ \\
    \cline{2-3}
    \multicolumn{1}{c}{} & \multicolumn{1}{c}{$f_{\cdot 1}$} & \multicolumn{1}{c}{$f_{\cdot 2}$} & $1$ \\
\end{tabular}
\label{tab:frequenz_oluanw}
\end{equation}


Die beobachteten relativen Häufigkeiten $f_{kj}$ werden mit den erwarteten Häufigkeiten
$\tilde{h}_{kj}$ verglichen. Die erwarteten Häufigkeiten werden wie in Gleichung
\ref{eq:erwartete_haeufigkeit} berechnet.

\begin{equation}
\tilde{f_{kj}} = f_{k\cdot} \cdot f_{\cdot j}
\label{eq:erwartete_haeufigkeit}
\end{equation}

Der $\chi^2$-Koeffizient wird berechnet, indem die quadrierten Differenzen
zwischen den beobachteten und erwarteten Häufigkeiten durch die erwarteten
Häufigkeiten geteilt und summiert werden, siehe Gleichung \ref{eq:chi_quadrat}.

\begin{equation}
\chi^2 = n \cdot \sum^K_{k = 1} \sum^2_{j = 1} \frac{(f_{kj}-\tilde{f_{kj}})^2}{\tilde{f_{kj}}}
\label{eq:chi_quadrat}
\end{equation}

Ist der $\chi^2$-Koeffizient groß, so ist die Abhängigkeit zwischen den beiden
Variablen Off-Label-Anwendung und $V$ stark. Ein kleiner $\chi^2$-Koeffizient
deutet darauf hin, dass die Off-Label-Anwendung $Anwendung^a$ und $V$ 
nicht unabhängig voneinander sind.

\subsubsection{Einschränkungen der Methodik}
\label{diskussion_oluanw}

Die gewählte Methodik bietet eine effiziente Möglichkeit zur Durchführung
deskriptiver Analyse der Off-Label-Anwendungen in verschiedenen Subgruppen.
Allerdings ist es wichtig, die Einschränkungen dieser Methodik zu
beachten.\bigskip

Ein wesentlicher Nachteil besteht darin, dass keine Inferenzstatistik angewendet
wird, was bedeutet, dass keine Aussagen über die Generalisierbarkeit der
Ergebnisse gemacht werden können. Die Ergebnisse sind ausschließlich auf die
untersuchten Daten beschränkt, und es können keine Schlussfolgerungen über ihre
Anwendbarkeit auf eine größere Population gezogen werden. Darüber hinaus können
aufgrund der geringen Datenmenge keine statistischen Tests wie der $\chi^2$-Test
durchgeführt werden, da die Anzahl der Fälle in den Zellen der
Häufigkeitstabelle zu klein ist.\bigskip

Selbst wenn ausreichend Daten vorhanden wären, müsste eine Korrektur für
multiple Tests durchgeführt werden, da mehrere Off-Label-Anwendungen untersucht
werden. Dies könnte zu einer erhöhten Wahrscheinlichkeit führen, einen Fehler 1.
Art zu begehen. Zudem ist zu beachten, dass die Anteile von Off-Label-Anwendungen
an allen Off-Label-Verordnungen nicht unabhängig voneinander sind, was weitere
Herausforderungen für die Interpretation der Ergebnisse mit sich bringt.\bigskip

Aus diesen Gründen sind die Ergebnisse tendenziell explorativer Natur und können
genutzt werden um Hypothesen für zukünftige Studien zu generieren.

\subsection{Ergebnisse}
\label{ergebnisse_oluanw}
Insgesamt wurden 26 verschiedene Off-Label-Anwendungen identifiziert, die
mindestens 20 Mal verschrieben wurden. Die Off-Label-Anwendungen wurden
bezüglich ihrem $\chi^2$-Koeffizienten sortiert, um die Unterschiede in der
Häufigkeit der Off-Label-Anwendungen in den verschiedenen Diagnosekategorien und
Einrichtungen systematisiert aufzuzeigen. In Tabelle \ref{table:oluanw_diag}
sehen wir für alle Off-Label-Anwendungen $Anwendung^a, a \in 1,\dots,A$ die
relativen Häufigkeiten $f_{\cdot 1}$ in den jeweiligen Diagnosekategorien, sowie
den $\chi^2$-Koeffizienten für die jeweilige Off-Label-Anwendung.

In Abbildung \ref{fig:diagnose_oluanw} sind diese relativen Häufigkeiten auch
visuell dargestellt, wobei die Off-Label-Anwendung nach den dazugehörigen
$\chi^2$-Koeffizienten sortiert sind. In den Plots wird ersichtlich, dass
beispielsweise die Off-Label-Anwendung Dexamethason\_Indikation Unterschiede in
der relativen Häufigkeit der Off-Label-Verordnungen in den verschiedenen
Diagnosekategorien hat, während die Off-Label-Anwendung Metamizol\_Modus, die im
letzten Plot dargestellt ist, weniger Unterschiede zwischen den Kategorien
aufweist.\bigskip

Mit dem gleichem Prinzip wurden die Abbildungen der Off-Label-Anwendungen im
Bezug auf die Einrichtungen erzeugt. Zunächst wurden die Off-Label-Anwendungen
unterteilt in ambulante und stationäre Einrichtungen und anschließend nach ihrem
$\chi^2$-Koeffizienten sortiert, wie es in der Abbildung
\ref{fig:einrichtung2_oluanw} aufgezeigt ist.\bigskip

Die Sortierung der Off-Label-Anwendungen gemäß Einrichtung2 wurde auch für
Einrichtung3, Einrichtung13 und Einrichtung16 verwendet, siehe Anhang A.2
Abbildungen \ref{fig:Einrichtung3_plot_oluanw},
\ref{fig:Einrichtung13_plot_oluanw} und \ref{fig:Einrichtung16_plot_oluanw}.

\begin{longtable}{lrrrr}
    \toprule
    Off-Label-Anwendung & internistisch & maligne & neurologisch & $\chi^2$-Koeffizient \\ 
    \midrule\addlinespace[2.5pt]
    Fentanyl\_Indikation & $4.56\%$ & $3.10\%$ & $3.37\%$ & $41.785$ \\ 
    Pantoprazol\_Indikation & $9.96\%$ & $10.33\%$ & $1.12\%$ & $39.982$ \\ 
    Metamizol\_Modus & $6.22\%$ & $6.44\%$ & $6.74\%$ & $28.862$ \\ 
    Mcp\_Modus & $0.83\%$ & $1.83\%$ & $1.12\%$ & $24.184$ \\ 
    Dexamethason\_Modus & $0.41\%$ & $1.75\%$ & $0.00\%$ & $17.443$ \\ 
    Hydromorphon\_Modus & $5.39\%$ & $4.29\%$ & $4.49\%$ & $14.762$ \\ 
    Haloperidol\_Applikationsweg & $4.98\%$ & $1.91\%$ & $4.49\%$ & $11.528$ \\ 
    Midazolam\_Modus & $5.39\%$ & $3.90\%$ & $7.87\%$ & $9.970$ \\ 
    Morphin\_Modus & $1.24\%$ & $1.83\%$ & $3.37\%$ & $6.742$ \\ 
    Midazolam\_Indikation & $7.47\%$ & $4.21\%$ & $8.99\%$ & $4.127$ \\ 
    Hydromorphon\_Dosis & $0.83\%$ & $2.23\%$ & $0.00\%$ & $3.999$ \\ 
    Haloperidol\_Indikation & $5.39\%$ & $4.13\%$ & $6.74\%$ & $3.945$ \\ 
    Natriumpicosulfat\_Dosis & $0.00\%$ & $1.91\%$ & $0.00\%$ & $3.335$ \\ 
    Torasemid\_Modus & $2.90\%$ & $1.75\%$ & $1.12\%$ & $2.845$ \\ 
    Morphin\_Indikation & $4.15\%$ & $4.53\%$ & $10.11\%$ & $2.475$ \\ 
    Butylscopolamin\_Modus & $2.49\%$ & $1.75\%$ & $5.62\%$ & $1.899$ \\ 
    Hydromorphon\_Indikation & $9.13\%$ & $6.92\%$ & $5.62\%$ & $1.027$ \\ 
    Dexamethason\_Indikation & $2.07\%$ & $8.43\%$ & $1.12\%$ & $0.684$ \\ 
    Mcp\_Dauer & $0.41\%$ & $3.18\%$ & $0.00\%$ & $0.443$ \\ 
    Buprenorphin\_Indikation & $0.83\%$ & $1.83\%$ & $1.12\%$ & $0.383$ \\ 
    Metamizol\_Applikationsweg & $3.73\%$ & $2.70\%$ & $3.37\%$ & $0.234$ \\ 
    Haloperidol\_Modus & $5.39\%$ & $2.94\%$ & $3.37\%$ & $0.001$ \\ 
    Midazolam\_Applikationsweg & $5.81\%$ & $2.70\%$ & $5.62\%$ & $0.001$ \\ 
    Mirtazapin\_Indikation & $2.90\%$ & $3.58\%$ & $1.12\%$ & $0.001$ \\ 
    Metoprolol\_Modus & $1.24\%$ & $1.59\%$ & $2.25\%$ & $0.000$ \\ 
    Butylscopolamin\_Indikation & $1.66\%$ & $1.83\%$ & $6.74\%$ & $0.000$ \\ 
    \bottomrule
    \caption{
    {Anteile der Off-Label-Anwendungen an den Off-Label-Verordnungen in den verschiedenen Diagnosekategorien.}
    }
    \label{table:oluanw_diag}
\end{longtable}

\begin{figure}[H]
    \centering
    \includegraphics[width=\textwidth]{../images/OLU_Anwendungen_Diagnose_plot.png}
    \caption{Anteile der Off-Label-Anwendungen an den Off-Label-Verordnungen in den verschiedenen Diagnosekategorien.}
    \label{fig:diagnose_oluanw}
\end{figure}

\begin{figure}[H]
    \centering
    \includegraphics[width=\textwidth]{../images/OLU_Anwendungen_Einrichtung2_plot.png}
    \caption{Anteile der Off-Label-Anwendungen an den Off-Label-Verordnungen in Einrichtung2.}
    \label{fig:einrichtung2_oluanw}
\end{figure}


